\documentclass[../main.tex]{subfiles}
\begin{document}

\chapter{Instruction Scheduling}
\lessoninfo{
  video={Lesson 7, videos 1--30},
  textbook={H\&P \cite{hennessy2017} §3.4--3.5; \cite{tomasulo1967}},
  keywords={out-of-order execution, Tomasulo's algorithm, instruction queue, reservation station, register alias table, tag, common data bus, issue, dispatch, write result, stale result, oldest first},
}

Lesson 6 showed that programs contain far more instruction-level parallelism
than an in-order processor can use, and that approaching it requires a
processor that renames registers and executes instructions out of order.
This lesson describes how such a processor decides, cycle by cycle, which
instructions to execute.  The method is Tomasulo's algorithm, which combines
register renaming with the detection of instructions whose operands are
available and is still the basis of today's out-of-order processors.  We
follow an instruction through its three steps, examine the choices that arise
when several instructions or results compete for the same hardware, trace a
complete example, and derive a shortcut for computing when each instruction
executes.

\section{Improving IPC}
\label{sec:l07-ipc}

Since the ILP of real programs is well above~1, a processor whose IPC is to
come close to it must deal with four kinds of dependences between
instructions.\source{Lesson 7, videos 1--2; H\&P §3.1.}  Each has its own
remedy:
\begin{description}
  \item[Control dependences] The instructions that follow a branch can enter
    the pipeline before the branch executes only if its outcome is guessed.
    Branch prediction (Lesson~4), and predication for hard-to-predict branches
    (Lesson~5), take care of them.
  \item[False data dependences] WAW and WAR dependences exist only because
    instructions reuse register names.  Register renaming (Lesson~6) removes
    them.
  \item[True data dependences] RAW dependences cannot be removed: an
    instruction has to wait until the values it reads have been
    computed.\caution{A dependence is a property of the program; a hazard is a
    property of the pipeline that has to resolve it.  H\&P §3.1.}
    What can be avoided is that instructions which do not depend on it wait as
    well.
  \item[Structural dependences] Instructions also compete for hardware:
    reservation stations, execution units and the bus introduced below.  The
    remedy is a wider processor (Lesson~6) with more of these resources, so
    that more instructions can be issued and executed in the same cycle.
\end{description}

\begin{definition}[Out-of-order execution]
  \label{def:l07-ooo}
  Executing each instruction as soon as the values it reads are available,
  even while instructions that precede it in program order are still waiting,
  is called \term{out-of-order execution}[アウトオブオーダ実行].
\end{definition}

Out-of-order execution is combined with renaming: without it, an instruction
could overtake only earlier instructions with which it has no WAW or WAR
dependence, and much of the ILP would be lost.  What remains is the
scheduling problem: in every cycle the processor must determine which
instructions have all their operands, and decide which of them to execute.

\section{Tomasulo's algorithm}
\label{sec:l07-tomasulo}

The scheduling method used in out-of-order processors was introduced in the
IBM System/360 Model~91.\source{Lesson 7, video 3; H\&P §3.4;
\cite{tomasulo1967}.}

\begin{definition}[Tomasulo's algorithm]
  \label{def:l07-tomasulo}
  The hardware algorithm published by Tomasulo in 1967, which renames
  registers and determines which instructions have their operands available
  so that each instruction executes as soon as its operands have been
  computed, is called \term{Tomasulo's algorithm}[Tomasulo のアルゴリズム].
\end{definition}

Current processors schedule instructions in essentially the same
way.\margin{Scale today: Skylake holds 97 instructions in its scheduler and
224 in its reorder buffer (Lesson~8).  Intel optimization manual.}  The
differences are listed in \cref{tab:l07-modern}.  The first two concern only
the scope of the algorithm; the third changes its mechanism.
Tomasulo's algorithm has no provision for exceptions: once instructions have
executed out of order, the registers need not correspond to any point in
the program at which an exception could be taken.  Supporting exceptions
requires an additional structure, the subject of
Lesson~8.\sidenote{The Model~91 (1967) ran at a \SI{60}{\nano\second} clock
and applied the algorithm to floating point only, where an add took 2 cycles
and a multiply~3; its interrupts were imprecise for exactly this reason.  The
idea then saw little use until the mid-1990s, when the Pentium~Pro and the
PowerPC~604 brought it back.}

\begin{table}[htbp]
  \centering\small
  \caption{Tomasulo's algorithm and modern out-of-order processors.}
  \label{tab:l07-modern}
  \begin{tabular}{@{}lll@{}}
    \toprule
    & Tomasulo's algorithm & Modern processors \\
    \midrule
    Instructions scheduled         & floating-point operations & all instructions \\
    Instructions in flight at once & a few                     & hundreds \\
    Exceptions                     & not supported             & supported (Lesson 8) \\
    \bottomrule
  \end{tabular}
\end{table}

\section{Organization}
\label{sec:l07-organization}

\Cref{fig:l07-tomasulo} shows the hardware that carries out Tomasulo's
algorithm.\source{Lesson 7, video 4; H\&P §3.4.}  The fetch stage appends
each fetched instruction to the \term{instruction queue}[命令キュー], a
first-in first-out buffer that keeps the fetched instructions in program
order until they can be moved on to the reservation stations.

\begin{figure}[htbp]
  \centering
  % Hardware organization of Tomasulo's algorithm: instruction queue, RAT and
% register file, reservation stations, execution units, common data bus, and
% the separate path of loads and stores (address adder, load and store
% buffers, memory).
\begin{tikzpicture}[x=1cm,y=1cm, font=\footnotesize\sffamily, >={Stealth[length=1.8mm]},
  blk/.style={draw, fill=ln-box, minimum height=0.6cm, inner sep=2pt, align=center},
  rs/.style={draw, fill=white, minimum width=1.5cm, minimum height=0.36cm, inner sep=0pt,
    font=\scriptsize\sffamily},
  qc/.style={draw, fill=white, minimum width=0.4cm, minimum height=0.5cm, inner sep=0pt},
  stp/.style={font=\scriptsize\bfseries\sffamily, text=ln-accent},
  bus/.style={line width=1.6pt, draw=ln-accent}]
  % The Japanese labels are wider than the English ones: shift the load and
  % store column to the right so that the "reservation stations" label fits
  % between the stations and the address adder (\lsd = 0 leaves EN unchanged).
  \iflangja{\def\lsd{0.8}}{\def\lsd{0}}
  % ---- fetch and instruction queue
  \node[blk, minimum width=1.1cm] (fetch) at (0.55,0) {\iflangja{フェッチ}{fetch}};
  \foreach \i in {0,...,5} { \node[qc] (q\i) at ({1.75+0.4*\i},0) {}; }
  \foreach \i in {2,...,5} { \node[qc, fill=ln-accent!15] at ({1.75+0.4*\i},0) {}; }
  \node[anchor=south, font=\scriptsize] at (2.75,0.28) {\iflangja{命令キュー}{instruction queue}};
  \draw[->] (fetch.east) -- (q0.west);
  % ---- issue: head of the queue to the stations, with the RAT
  \draw[->] (q5.east) -- (4.6,0) -- (4.6,-0.75);
  \node[stp, anchor=west] at (4.65,-0.3) {\iflangja{発行}{issue}};
  \draw (0.55,-0.75) -- ({9.3+\lsd},-0.75);
  \fill (4.6,-0.75) circle (1.3pt);
  % ---- RAT and register file
  \node[blk, fill=ln-accent!15, minimum width=1.1cm] (rat) at (0.55,-1.45) {RAT};
  \node[blk, minimum width=1.1cm, font=\scriptsize\sffamily] (regs) at (0.55,-2.35) {\iflangja{レジスタ\\ファイル}{register\\file}};
  \draw[<->] (rat.north) -- (0.55,-0.75);
  \draw (rat.south) -- (regs.north);
  % ---- reservation stations
  \foreach \n/\y in {1/-1.25,2/-1.61,3/-1.97} { \node[rs] (ad\n) at (2.9,\y) {AD\n}; }
  \foreach \n/\y in {1/-1.25,2/-1.61} { \node[rs] (ml\n) at (6.1,\y) {ML\n}; }
  \draw[->] (2.9,-0.75) -- (ad1.north);
  \draw[->] (6.1,-0.75) -- (ml1.north);
  \node[anchor=west, font=\scriptsize, text=ln-muted, align=left] at (6.95,-1.43)
    {\iflangja{リザベーション\\ステーション}{reservation\\stations}};
  % ---- load/store path: address adder, load and store buffers
  \node[blk, minimum width=1.6cm, minimum height=0.5cm] (aadd) at ({9.3+\lsd},-1.35) {\iflangja{アドレス加算器}{address adder}};
  \draw[->] ({9.3+\lsd},-0.75) -- (aadd.north);
  \node[anchor=south east, font=\scriptsize, text=ln-muted] at ({9.3+\lsd},-0.75) {LW, SW};
  \node[blk, minimum width=1.1cm, font=\scriptsize\sffamily] (sb) at ({8.7+\lsd},-2.3) {\iflangja{ストア\\バッファ}{store\\buffers}};
  \node[blk, minimum width=1.1cm, font=\scriptsize\sffamily] (lb) at ({9.9+\lsd},-2.3) {\iflangja{ロード\\バッファ}{load\\buffers}};
  \draw (aadd.south) -- ({9.3+\lsd},-1.8);
  \draw[->] ({9.3+\lsd},-1.8) -| (sb.north);
  \draw[->] ({9.3+\lsd},-1.8) -| (lb.north);
  % ---- execution units
  \node[blk, minimum width=1.5cm] (adder) at (2.9,-3.3) {\iflangja{加減算器}{adder}};
  \node[blk, minimum width=1.5cm] (mult) at (6.1,-3.3) {\iflangja{乗除算器}{multiplier}};
  \node[blk, minimum width=2.3cm] (mem) at ({9.3+\lsd},-3.3) {\iflangja{メモリ}{memory}};
  \draw[->] (ad3.south) -- (adder.north);
  \draw[->] (ml2.south) -- (mult.north);
  \draw[->] (sb.south) -- (sb.south |- mem.north);
  \draw[->] (lb.south) -- (lb.south |- mem.north);
  \node[stp, anchor=west] at (2.95,-2.6) {\iflangja{ディスパッチ}{dispatch}};
  \node[stp, anchor=west] at (6.15,-2.6) {\iflangja{ディスパッチ}{dispatch}};
  % ---- common data bus and its destinations
  \draw[bus] (0.55,-4.3) -- ({10.6+\lsd},-4.3);
  \node[anchor=north east, font=\scriptsize\bfseries\sffamily, text=ln-accent] at ({10.6+\lsd},-4.38)
    {CDB\enskip(\iflangja{結果書き込み}{write result})};
  \draw[->] (adder.south) -- (2.9,-4.3);
  \draw[->] (mult.south) -- (6.1,-4.3);
  \draw[->] (mem.south) -- ({9.3+\lsd},-4.3);
  \draw[->, draw=ln-accent] (0.55,-4.3) -- (regs.south);
  \draw[->, draw=ln-accent] (1.65,-4.3) -- (1.65,-1.61) -- ([xshift=-0.75cm]ad2.center);
  \draw[->, draw=ln-accent] (4.9,-4.3) -- (4.9,-1.43) -- ([xshift=-0.75cm]6.1,-1.43);
  \draw[->, draw=ln-accent] (1.65,-2.9) -- (1.25,-2.9) |- (rat.east);
  \draw[->, draw=ln-accent] (7.8,-4.3) -- (7.8,-2.3) -- (sb.west);
\end{tikzpicture}

  \caption{Organization of a processor with Tomasulo's algorithm.  Issue moves
    the instruction at the head of the queue into a reservation station, using
    the RAT and the register file; dispatch sends a ready instruction to its
    execution unit; results return over the common data bus to the
    reservation stations, the RAT, the register file and the store buffers.
    Loads and stores take the path through the address adder and the load
    and store buffers.}
  \label{fig:l07-tomasulo}
\end{figure}

\begin{definition}[Reservation station]
  \label{def:l07-rs}
  A \term{reservation station}[リザベーションステーション]%
  \index{RS|see{reservation station}} (RS) holds one instruction from the
  moment it leaves the instruction queue until its result has been
  broadcast.  It records whether it is busy; the operation; for each operand
  either its value or, if the value has not been computed yet, the tag of the
  value; and whether the instruction has already been dispatched.
\end{definition}

Reservation stations are grouped by the execution unit that serves them.  The
examples of this lesson use three stations AD1, AD2 and AD3 for additions and
subtractions, served by one adder, and two stations ML1 and ML2 for
multiplications and divisions, served by one multiplier.\margin{The 360/91
had exactly this arrangement: three reservation stations on its floating-point
adder and two on its multiplier--divider \cite{tomasulo1967}.}  The name of a
station is the \emph{tag} under which the result of its instruction will be
broadcast.

The register alias table of Lesson~6 has one entry per register.  An empty
entry means that the register file holds the current value of the register;
otherwise the entry holds the tag of the station whose instruction will
produce that value.  Tags thus play the role of the physical registers of
Lesson~6: every result receives a new name, and an instruction that reads a
register refers to the result of its closest preceding writer.

\begin{definition}[Common data bus]
  \label{def:l07-cdb}
  The \term{common data bus}[共通データバス]%
  \index{CDB|see{common data bus}} (CDB) carries each result computed by an
  execution unit, together with its tag, to all reservation stations, to the
  RAT and to the register file.
\end{definition}

Loads and stores use a separate path.  An adder computes the memory address;
a load then waits in a load buffer, which holds the address, and a store in a
store buffer, which holds the address and the value to be stored, possibly
received from the CDB.  The value read by a load is placed on the CDB like any
other result (\cref{sec:l07-load-store}).

Every instruction passes through three steps, marked in
\cref{fig:l07-tomasulo}:
\begin{description}
  \item[Issue] The instruction at the head of the instruction queue is placed
    in a free reservation station, with its operands and its destination
    resolved through the RAT.  This step is called
    \term{issue}[発行].
  \item[Dispatch] An instruction whose operands are all available is sent
    from its reservation station to an execution unit; this is its
    \term{dispatch}[ディスパッチ].
  \item[Write result] The result is placed on the CDB, from which it reaches
    the waiting reservation stations, the RAT and the register file.  This
    step is called \term{write result}[結果書き込み], or
    broadcast.\index{broadcast|see{write result}}
\end{description}

\section{Issue}
\label{sec:l07-issue}

Issue takes the instruction at the head of the instruction queue and performs
the following steps.\source{Lesson 7, videos 5--7; H\&P §3.5.}  In the
examples of this lesson one instruction is issued per cycle; a wider
processor issues several per cycle, still in program order.
\begin{enumerate}
  \item Take the next instruction, in program order, from the head of the
    queue.
  \item Determine where each operand comes from.  If the RAT entry of the
    operand register is empty, the value is in the register file; otherwise
    the entry names the tag of the value.
  \item Find a free reservation station for the operation.  If every station
    of that kind is busy, the instruction is not issued in this cycle; it
    stays at the head of the queue, and the instructions behind it wait as
    well.
  \item Place the instruction in the station: the operation, and for each
    operand the value from the register file or the tag from the RAT.
  \item Set the RAT entry of the destination register to the name of the
    station.
\end{enumerate}

Instructions are issued strictly in program order.  The RAT then describes the
register values seen at the position of the instruction being issued, so that
each operand refers to its closest preceding writer.  As in renaming, the
operands are looked up before the destination entry is
overwritten.\caution{Issue never waits for operands, only for a free
reservation station.}

\begin{example}
  \label{ex:l07-program}
  The following program is the running example of this lesson:
  \begin{lstlisting}[language={[x86masm]Assembler}, numbers=none]
I1:  MUL R1, R2, R3     ; R1 <- R2 * R3
I2:  ADD R4, R1, R5     ; R4 <- R1 + R5
I3:  SUB R6, R2, R5     ; R6 <- R2 - R5
I4:  ADD R1, R5, R3     ; R1 <- R5 + R3
I5:  DIV R2, R6, R3     ; R2 <- R6 / R3
I6:  MUL R3, R4, R1     ; R3 <- R4 * R1
I7:  SUB R5, R3, R2     ; R5 <- R3 - R2
\end{lstlisting}
  Initially the RAT is empty and the register file holds R1 $=1$,
  R2 $=14$, R3 $=4$, R4 $=3$, R5 $=6$ and R6 $=7$.  In cycle~1, I1 is issued
  into ML1 with the operand values 14 and 4, and the RAT entry of R1 becomes
  ML1.  In cycle~2, I2 is issued into AD1 (\cref{fig:l07-issue}).  The RAT
  entry of R1 holds ML1, so the first operand is the tag ML1; the entry of R5
  is empty, so the second operand is the value~6.  The RAT entry of R4 becomes
  AD1.  I3 and I4 follow in cycles 3 and~4.  I4 writes R1 again and replaces
  ML1 with AD3 in the entry of R1: every instruction issued from now on that
  reads R1 uses the result of I4, not that of I1.
\end{example}

\begin{figure}[htbp]
  \centering
  % Issuing I2 (ADD R4, R1, R5) in cycle 2 of the running example: the RAT
% entry of R1 holds the tag ML1, the entry of R5 is empty so the value comes
% from the register file, and the RAT entry of R4 is set to AD1.
\begin{tikzpicture}[x=1cm,y=1cm, font=\footnotesize\sffamily, >={Stealth[length=1.8mm]},
  cl/.style={draw, fill=white, minimum width=0.9cm, minimum height=0.45cm, inner sep=0pt},
  nm/.style={cl, fill=ln-box, minimum width=0.8cm},
  new/.style={cl, fill=ln-accent!15},
  tg/.style={font=\footnotesize\itshape\sffamily},
  lab/.style={font=\scriptsize\sffamily, text=ln-muted},
  arr/.style={->, thick, ln-accent}]
  % ---- register file and RAT, one column per register
  \node[anchor=west] at (-0.8,1.05) {I2:\enskip\texttt{ADD R4, R1, R5}};
  \foreach \r/\x in {1/1.45,2/2.35,3/3.25,4/4.15,5/5.05,6/5.95} {
    \node[lab] at (\x,0.4) {R\r};
  }
  \node[lab] at (1.45,0.66) {src$_1$};
  \node[lab] at (4.15,0.66) {dst};
  \node[lab] at (5.05,0.66) {src$_2$};
  \node[lab, anchor=east] at (0.95,0) {\iflangja{レジスタ}{register file}};
  \node[lab, anchor=east] at (0.95,-0.45) {RAT};
  \foreach \v/\x in {1/1.45,14/2.35,4/3.25,3/4.15,6/5.05,7/5.95} {
    \node[cl] at (\x,0) {\v};
  }
  \node[cl, tg] (r1) at (1.45,-0.45) {ML1};
  \foreach \x in {2.35,3.25,5.05,5.95} { \node[cl] at (\x,-0.45) {}; }
  \node[new, tg] (r4) at (4.15,-0.45) {AD1};
  \node[cl, draw=none, fill=none] (r5) at (5.05,-0.45) {};
  % ---- reservation stations: adder group left, multiplier group right
  \foreach \n/\y in {1/-1.6,2/-2.05,3/-2.5} {
    \node[nm] (ad\n) at (0.4,\y) {AD\n};
    \node[cl, minimum width=0.8cm] (ado\n) at (1.2,\y) {};
    \node[cl] (ada\n) at (2.05,\y) {};
    \node[cl] (adb\n) at (2.95,\y) {};
  }
  \foreach \n/\y in {1/-1.6,2/-2.05} {
    \node[nm] (ml\n) at (4.4,\y) {ML\n};
    \node[cl, minimum width=0.8cm] (mlo\n) at (5.2,\y) {};
    \node[cl] (mla\n) at (6.05,\y) {};
    \node[cl] (mlb\n) at (6.95,\y) {};
  }
  \node[new, minimum width=0.8cm] at (1.2,-1.6) {ADD};
  \node[new, tg] (opa) at (2.05,-1.6) {ML1};
  \node[new] (opb) at (2.95,-1.6) {6};
  \node at (5.2,-1.6) {MUL};
  \node at (6.05,-1.6) {14};
  \node at (6.95,-1.6) {4};
  \node[lab] at (1.7,-2.95) {\iflangja{加減算用}{for the adder}};
  \node[lab] at (5.7,-2.5) {\iflangja{乗除算用}{for the multiplier}};
  % ---- operand lookups
  \draw[arr] (r1.south) -- (opa.north);
  \draw[arr] (5.05,-0.225) -- (r5.south) -- (opb.north);
\end{tikzpicture}

  \caption{Issue of I2 in cycle~2 of \cref{ex:l07-program}.  The RAT entry of
    R1 yields the tag ML1 (tags in italics), the empty entry of R5 lets the
    value~6 be taken from the register file, and the entry of R4 is set to
    AD1.  Shaded cells are written in this step.}
  \label{fig:l07-issue}
\end{figure}

\section{Dispatch}
\label{sec:l07-dispatch}

An instruction waits in its reservation station until all its operands are
values.\source{Lesson 7, video 8; H\&P §3.5.}  Whenever a result appears
on the CDB, every reservation station compares the tags of its operands with
the tag on the bus and, where they match, replaces the tag with the value
from the bus; this is called \emph{capturing} the value.  A station whose
operands are all values holds a \emph{ready} instruction.\caution{Other texts
swap the names: \emph{dispatch} for placing an instruction in a station and
\emph{issue} for sending it to a unit.  H\&P's steps are issue, execute and
write result.}  When the execution
unit that serves the station is free, a ready instruction is dispatched: its
operation and operand values are sent to the unit, which starts executing it.
The station remains busy while the instruction executes, because its name is
the tag under which the result will be broadcast.

\subsection{More than one ready instruction}

Several instructions may be ready when an execution unit becomes free, and
one of them must be chosen.\source{Lesson 7, videos 9--10.}  Every choice is
correct, since a ready instruction has all the values it needs; the choice
affects only performance and hardware cost.\sidenote{Wakeup, the comparison of
the broadcast tag against every waiting operand, and select, the choice of one
ready instruction per unit, must together fit in one cycle; otherwise
dependent instructions cannot execute back to back.  The cost of that loop
grows with the number of stations, which is why schedulers hold about a
hundred instructions where reorder buffers hold hundreds
\cite{palacharla1997}.}
\begin{description}
  \item[Oldest first]\index{oldest first} Dispatch the instruction that was
    issued first.  Other things being equal, more of the instructions in
    flight were issued after an older instruction and may be waiting for its
    result, so this approximates the next policy while requiring only that
    the order of issue be recorded.
  \item[Most dependents first] Dispatch the instruction whose result is
    needed by the largest number of waiting instructions, so that as many
    instructions as possible become ready soon.  This promises the best performance, but
    finding that instruction means comparing the tag of every candidate with
    the operands of all reservation stations, which is too slow and costly.
  \item[Random] The simplest selection logic.  It may keep an instruction on
    which many others depend waiting behind unimportant ones, so performance
    suffers.
\end{description}
Oldest first is a good compromise between the two others and is the rule used
in the examples of this lesson.  In the running example, both I2 and I4 are ready when the adder becomes free in cycle~7
(\cref{sec:l07-example}); oldest first dispatches I2, although I4 has been
ready two cycles longer.

\section{Write result}
\label{sec:l07-write}

When an execution unit has finished an instruction, the write-result step
broadcasts its result.\source{Lesson 7, video 11; H\&P §3.5.}
\begin{enumerate}
  \item The result and the tag of the instruction's reservation station are
    placed on the CDB.
  \item Every operand that holds this tag captures the value
    (\cref{sec:l07-dispatch}).
  \item If the RAT entry of a register holds the tag, the value is written
    into that register in the register file and the RAT entry is cleared, so
    that instructions issued later read the value from the register file.
  \item The reservation station is freed.
\end{enumerate}
The station need not record the destination register, because the RAT entry
that holds the tag identifies it.  \Cref{fig:l07-broadcast} shows the
broadcast of the result of I2 in the running example.

\begin{figure}[htbp]
  \centering
  % Write result of I2 in cycle 9 of the running example: tag AD1 and value 62
% on the CDB.  ML1 captures the value, the RAT entry of R4 matches, so R4 is
% written and its entry cleared, and AD1 is freed.
\begin{tikzpicture}[x=1cm,y=1cm, font=\footnotesize\sffamily, >={Stealth[length=1.8mm]},
  cl/.style={draw, fill=white, minimum width=0.9cm, minimum height=0.45cm, inner sep=0pt},
  nm/.style={cl, fill=ln-box, minimum width=0.8cm},
  new/.style={cl, fill=ln-accent!15},
  tg/.style={font=\footnotesize\itshape\sffamily},
  lab/.style={font=\scriptsize\sffamily, text=ln-muted},
  arr/.style={->, thick, ln-accent}]
  % ---- register file and RAT
  \foreach \r/\x in {1/1.45,2/2.35,3/3.25,4/4.15,5/5.05,6/5.95} {
    \node[lab] at (\x,0.4) {R\r};
  }
  \node[lab, anchor=east] at (0.95,0) {\iflangja{レジスタファイル}{register file}};
  \node[lab, anchor=east] at (0.95,-0.45) {RAT};
  \foreach \v/\x in {1/1.45,14/2.35,4/3.25,6/5.05,8/5.95} {
    \node[cl] at (\x,0) {\v};
  }
  \node[new] (g4) at (4.15,0) {62};
  \foreach \t/\x in {AD3/1.45,ML2/2.35,ML1/3.25,AD2/5.05} {
    \node[cl, tg] at (\x,-0.45) {\t};
  }
  \node[cl] at (5.95,-0.45) {};
  \node[new] (r4) at (4.15,-0.45) {};
  % ---- common data bus
  \draw[line width=1.6pt, ln-accent] (-0.1,-1.05) -- (7.5,-1.05);
  \node[anchor=west, font=\scriptsize\bfseries\sffamily, text=ln-accent, align=left]
    at (7.55,-1.05) {CDB\\\iflangja{タグ}{tag} \textit{AD1}, \iflangja{値}{value} 62};
  % ---- reservation stations
  \foreach \n/\y in {1/-1.6,2/-2.05,3/-2.5} {
    \node[nm] (ad\n) at (0.4,\y) {AD\n};
    \node[cl, minimum width=0.8cm] at (1.2,\y) {};
    \node[cl] at (2.05,\y) {};
    \node[cl] at (2.95,\y) {};
  }
  \foreach \n/\y in {1/-1.6,2/-2.05} {
    \node[nm] (ml\n) at (4.4,\y) {ML\n};
    \node[cl, minimum width=0.8cm] at (5.2,\y) {};
    \node[cl] at (6.05,\y) {};
    \node[cl] at (6.95,\y) {};
  }
  % AD1 (I2) is freed
  \node[lab] at (2.05,-1.6) {\iflangja{解放}{freed}};
  % AD2 (I7), AD3 (I4), ML2 (I5): no operand waits for AD1
  \node at (1.2,-2.05) {SUB};  \node[tg] at (2.05,-2.05) {ML1}; \node[tg] at (2.95,-2.05) {ML2};
  \node at (1.2,-2.5) {ADD};   \node at (2.05,-2.5) {6};         \node at (2.95,-2.5) {4};
  \node at (5.2,-2.05) {DIV};  \node at (6.05,-2.05) {8};         \node at (6.95,-2.05) {4};
  % ML1 (I6) captures the value in its first operand
  \node at (5.2,-1.6) {MUL};
  \node[new] (cap) at (6.05,-1.6) {62};
  \node[tg] at (6.95,-1.6) {AD3};
  % ---- effects of the broadcast
  \draw[arr] (4.15,-1.05) -- (g4.south);
  \draw[arr] (6.05,-1.05) -- (cap.north);
  \fill[ln-accent] (4.15,-1.05) circle (1.5pt) (6.05,-1.05) circle (1.5pt);
  \node[lab, anchor=west] at (4.2,-0.83) {\iflangja{一致}{match}};
\end{tikzpicture}

  \caption{Write result of I2 in cycle~9 of the running example.  ML1 (I6)
    captures 62 in place of the tag AD1; the RAT entry of R4 holds AD1, so
    R4 becomes 62 and the entry is cleared; AD1 is freed.  No other operand
    or RAT entry holds AD1.}
  \label{fig:l07-broadcast}
\end{figure}

\subsection{More than one result in a cycle}
\label{sec:l07-cdb-conflict}

Two execution units may finish in the same cycle, but one bus carries one
result at a time.\source{Lesson 7, video 12.}  One solution is a separate
bus for each unit.  Every operand of every reservation station must then
compare its tag with each bus and select the value from the matching one, and
the register file needs a write port per bus; the comparison logic, which is
already large, grows in proportion to the number of buses.  The alternative
keeps a single bus and lets one result wait for the next cycle.  The rule
used here gives priority to the result of the \emph{slower} unit: its
instruction was dispatched longer ago, so it tends to be the older
instruction, and more instructions are likely to be waiting for its result.
In the examples we assume that a unit whose result is waiting for the bus
cannot start another instruction.

\subsection{Stale results}
\label{sec:l07-stale}

By the time a result is broadcast, a later instruction that writes the same
register may already have been issued.\source{Lesson 7, video 13; H\&P §3.5.}
The RAT entry of the register then holds the tag of the later instruction,
not the tag on the bus, and step~3 has no effect: neither the register file
nor the RAT changes.  The result is \emph{stale}\index{stale result} only for
the register, and ignoring
it there is correct.  After both instructions the register must hold the later
result, which arrives when the later instruction broadcasts.  The earlier
result is still needed, but only by the instructions issued between the two
writers; they hold its tag and capture the value in step~2.  Instructions
issued after the later writer refer to the later tag.

\begin{example}
  \label{ex:l07-stale}
  In the running example, I1 and I4 both write R1.  When I1 broadcasts 56
  under the tag ML1 in cycle~6, the RAT entry of R1 already holds AD3, because
  I4 was issued in cycle~4.  R1 in the register file keeps its old value~1
  and the entry keeps AD3, while AD1, which holds I2, captures~56.  R1 becomes
  10 only in cycle~11, when I4 broadcasts under AD3.
\end{example}

\begin{keypoint}
  A broadcast delivers its value to every operand that waits for its tag, but
  it updates the register file and the RAT only if the RAT entry of the
  destination still holds that tag.
\end{keypoint}

\section{The steps in one cycle}
\label{sec:l07-cycle}

For a single instruction, Tomasulo's algorithm is a sequence: issue, waiting
and capturing in a reservation station, dispatch, execution, write
result.\source{Lesson 7, videos 14--19; H\&P §3.5.}  For the processor as a
whole, all steps take place in every cycle, each for different instructions:
in the same cycle one instruction may be issued, several stations may capture
a broadcast value, ready instructions may be dispatched to free units, and one
result is broadcast.

Whether two steps of the \emph{same} instruction can take place in the same
cycle is a matter of circuit design.  In the simplest design every step
examines the structures as they are at the start of the cycle, and all
updates take effect at its end.  Then:
\begin{itemize}
  \item An instruction issued in cycle $C$ is in its reservation station only
    from cycle $C+1$, so it can be dispatched in $C+1$ at the earliest.
  \item An instruction that captures its last missing operand in cycle $C$ is
    ready from cycle $C+1$ and can be dispatched in $C+1$ at the earliest.
  \item A station freed by a broadcast in cycle $C$ can receive a new
    instruction in cycle $C+1$.
\end{itemize}
The delay between capture and dispatch can be removed with additional logic,
a dispatch that also considers operands arriving on the bus in the current
cycle, and a station freed by a broadcast could similarly be reused in the
same cycle.  The additional logic lies in the same cycle as the step it
extends and therefore lengthens the clock cycle.  An instruction, however, is
not dispatched in the cycle in which it is issued: while issue writes it into
its station, dispatch does not yet see it there.

Two interactions between issue and write result in the same cycle need care.
\begin{itemize}
  \item Both can change the RAT entry of the same register: the broadcast of
    an older writer clears it, and the issue of a newer writer sets it to a
    new tag.  The issue must prevail, because it belongs to the later
    instruction in program order; equivalently, the write result is applied
    first and the issue second.  Otherwise later instructions would read the
    register file and miss the result of the newly issued instruction.
  \item An instruction may be issued in the cycle in which one of its operands
    is broadcast.  In the simplest design it finds the tag in the RAT, since
    the entry is cleared only at the end of the cycle, but enters its station
    after the value has passed on the bus; issue must then compare operand
    tags with the bus and take the value directly, or the instruction would
    wait for a broadcast that never comes.  If the updates of the broadcast
    are visible to the issue in the same cycle, the issue simply reads the
    new value from the register file.
\end{itemize}

\begin{keypoint}
  In the simplest design an instruction advances by at most one step per
  cycle: issued in cycle $C$, it is dispatched in $C+1$ at the earliest, and a
  result broadcast in cycle $W$ lets the instructions waiting for it be
  dispatched in $W+1$ at the earliest.
\end{keypoint}

\section{Load and store instructions}
\label{sec:l07-load-store}

Loads and stores create dependences through memory that are analogous to the
dependences through registers.\source{Lesson 7, video 20; H\&P §3.5.}
\begin{lstlisting}[language={[x86masm]Assembler}, numbers=none]
SW  R1, 0(R2)      ; MEM[R2] <- R1
LW  R3, 8(R4)      ; R3 <- MEM[R4 + 8]
\end{lstlisting}
If R4${}+8$ equals R2, the load reads the value just stored, a RAW
dependence through memory.  Two stores to the same address have a WAW
dependence, and a store to an address that an earlier load reads has a WAR
dependence.  The RAT cannot remove the false ones, because it renames
registers, not memory locations, and whether two memory accesses conflict is
known only once both addresses have been computed.\margin{Skylake can have 72
loads and 56 stores in flight.  Intel optimization manual.}

Tomasulo's algorithm therefore performs loads and stores in program
order.\caution{\enquote{In program order} refers to the memory accesses.  The
addresses themselves are computed as soon as their operands are ready, and
therefore out of order.}  They are issued like other instructions, with their operands resolved through
the RAT, and the address adder computes the address.  A load then waits in a
load buffer and a store in a store buffer (\cref{fig:l07-tomasulo}); a store
whose value has not been computed yet captures it from the CDB.  The
memory accesses are made from these buffers in the order in which the loads
and stores were issued.  A load broadcasts the value it reads on
the CDB under its own tag, exactly like an arithmetic result.  Modern
processors also execute loads and stores out of order, which requires the
processor to detect and handle the memory dependences; Lesson~9 covers this.\margin{Lesson~9 lets a
load overtake a store whose address is not yet known, and checks afterwards
whether the two addresses overlapped.}

\section{A complete example}
\label{sec:l07-example}

We now execute the program of \cref{ex:l07-program} to the
end.\source{Lesson 7, videos 21--26; H\&P §3.5.}  The assumptions are those of
the simplest design of \cref{sec:l07-cycle}, together with the following:
\begin{itemize}
  \item At most one instruction is issued per cycle, into the free station of
    the required kind with the lowest number.
  \item The adder executes ADD and SUB in 2 cycles; the multiplier executes
    MUL and DIV in 4 cycles.\margin{For scale: on Skylake a floating-point add
    or multiply takes 4 cycles, a double-precision divide 13--14 and an
    integer add~1.  Agner Fog, \emph{Instruction tables}.}  An instruction
    dispatched in cycle $D$ with an execution time of $k$ cycles executes in
    cycles $D$ to $D+k-1$ and broadcasts in cycle $D+k$, in which its unit
    can already accept the next instruction.\margin{The units here are not pipelined: each is busy in
    cycles $D$ to $D+k-1$.  Real adders and multipliers accept a new operation
    every cycle.}
  \item There is one CDB, and the multiplier has priority over the adder
    (\cref{sec:l07-cdb-conflict}); a unit whose result waits for the bus
    accepts no new instruction.  Ready instructions are dispatched oldest
    first.
\end{itemize}
\Cref{tab:l07-trace} lists what happens in each cycle, and
\cref{tab:l07-state} shows the reservation stations and the RAT after cycles
5, 6 and~7.

\begin{table}[htbp]
  \centering\small
  \caption{Events in each cycle of the complete example.  Captures are listed
    with the broadcast that causes them.}
  \label{tab:l07-trace}
  \begin{tabular}{@{}c>{\raggedright\arraybackslash}p{0.14\linewidth}>{\raggedright\arraybackslash}p{0.12\linewidth}>{\raggedright\arraybackslash}p{0.55\linewidth}@{}}
    \toprule
    Cycle & Issue & Dispatch & Write result and remarks \\
    \midrule
    1     & I1 $\to$ ML1 &    & \\
    2     & I2 $\to$ AD1 & I1 & \\
    3     & I3 $\to$ AD2 &    & I2 waits for ML1 \\
    4     & I4 $\to$ AD3 & I3 & RAT entry of R1 becomes AD3 \\
    5     & I5 $\to$ ML2 &    & all stations busy; I4 ready, adder busy \\
    6     &              &    & I1: ML1 $=56$, stale for R1; AD1 captures;
                                I3 finished, waits for the bus;
                                no free station for I6 \\
    7     & I6 $\to$ ML1 & I2 & I3: AD2 $=8$, R6 $\leftarrow 8$; ML2 captures;
                                I2 and I4 ready, I2 older \\
    8     & I7 $\to$ AD2 & I5 & \\
    9     &              & I4 & I2: AD1 $=62$, R4 $\leftarrow 62$; ML1 captures \\
    10    &              &    & \\
    11    &              &    & I4: AD3 $=10$, R1 $\leftarrow 10$; ML1 captures \\
    12    &              & I6 & I5: ML2 $=2$, R2 $\leftarrow 2$; AD2 captures \\
    13--15 &             &    & I6 executes \\
    16    &              &    & I6: ML1 $=620$, R3 $\leftarrow 620$; AD2 captures \\
    17    &              & I7 & \\
    18    &              &    & \\
    19    &              &    & I7: AD2 $=618$, R5 $\leftarrow 618$ \\
    \bottomrule
  \end{tabular}
\end{table}

\begin{table}[htbp]
  \centering\small
  \caption{Reservation stations and RAT at the end of cycles 5, 6 and~7.
    Tags are in italics; \enquote{exec} marks an instruction that has been
    dispatched.  An empty RAT entry is shown with the register-file value.}
  \label{tab:l07-state}
  \begin{tabular}{@{}llll@{}}
    \toprule
    & After cycle 5 & After cycle 6 & After cycle 7 \\
    \midrule
    AD1 & I2: \textit{ML1}, 6    & I2: 56, 6              & I2: 56, 6 (exec) \\
    AD2 & I3: 14, 6 (exec)       & I3: 14, 6 (exec)       & free \\
    AD3 & I4: 6, 4               & I4: 6, 4               & I4: 6, 4 \\
    ML1 & I1: 14, 4 (exec)       & free                   & I6: \textit{AD1}, \textit{AD3} \\
    ML2 & I5: \textit{AD2}, 4    & I5: \textit{AD2}, 4    & I5: 8, 4 \\
    \midrule
    R1  & \textit{AD3} & \textit{AD3} & \textit{AD3} \\
    R2  & \textit{ML2} & \textit{ML2} & \textit{ML2} \\
    R3  & 4            & 4            & \textit{ML1} \\
    R4  & \textit{AD1} & \textit{AD1} & \textit{AD1} \\
    R5  & 6            & 6            & 6 \\
    R6  & \textit{AD2} & \textit{AD2} & 8 \\
    \bottomrule
  \end{tabular}
\end{table}

\paragraph{Cycles 1--5.}  One instruction is issued per cycle.  I1 and I3
have all their operands at issue and are dispatched in the next cycle, I1 to
the multiplier and I3 to the adder.  I4 is ready from cycle~5, but the adder
is busy with I3.  I2 waits for the tag ML1 and I5 for AD2.  After cycle~5
every station is busy.

\paragraph{Cycle 6.}  Both units finished in cycle~5.  The bus takes the
result of the slower multiplier, 56 under ML1, which is stale for R1
(\cref{ex:l07-stale}); AD1 captures it and ML1 is freed.  The result of I3
waits in the adder.  I6 is not issued, because at the start of the cycle
both multiplier stations are busy.

\paragraph{Cycle 7.}  I3 broadcasts 8 under AD2.  The RAT entry of R6 holds
AD2, so R6 becomes~8 and the entry is cleared; ML2 captures the value for I5,
and AD2 is freed.  The adder can take a new instruction, and both I2 and I4
are ready; oldest first dispatches I2.  I6 is issued into ML1, freed in the
previous cycle, and receives two tags, AD1 and AD3.

\paragraph{Cycles 8--19.}  I7 is issued into AD2 and waits for ML1 and ML2.
I5 starts on the multiplier in cycle~8, and I4 on the adder in cycle~9, when I2
releases it.  I6 needs the results of I2 (cycle~9) and I4 (cycle~11) and the
multiplier, which I5 releases in cycle~12; I7 needs the result of I6,
broadcast in cycle~16.  After the last broadcast, in cycle~19, the register
file holds R1 $=10$, R2 $=2$, R3 $=620$, R4 $=62$, R5 $=618$ and R6
$=8$,\margin{IPC here is $7/19 \approx 0.37$: even with unlimited stations and
units the chain I1$\to$I2$\to$I6$\to$I7 would need 17 cycles.}
exactly the values produced by executing the program in order.

\begin{keypoint}
  The instructions of the example broadcast in the order I1, I3, I2, I4, I5,
  I6, I7, and each executes as soon as its operands, a free station and a free
  unit allow.  Because every operand is resolved through a tag at issue and
  stale results never reach the register file, the final register values are
  those of sequential execution.
\end{keypoint}

\section{Timing without tracing the structures}
\label{sec:l07-timing}

Tracing every structure cycle by cycle is laborious.  When only the timing
matters, it suffices to determine the cycles of issue, dispatch and write
result of each instruction.\source{Lesson 7, videos 27--29; H\&P §3.5.}  Let
$I_n$, $D_n$ and $W_n$ denote these cycles for instruction $n$, and $k_n$ its
execution time.  Under the assumptions of \cref{sec:l07-example}:
\begin{description}
  \item[Issue] $I_n$ is the first cycle after $I_{n-1}$ in which a station
    of the required kind is free.  A station freed by the write result of
    instruction $m$ is free from cycle $W_m+1$.
  \item[Dispatch] $D_n$ is the first cycle that satisfies
    \begin{equation}
      D_n \;\ge\; 1 + \max\Bigl(I_n,\; \max_{p \in P(n)} W_p\Bigr) ,
      \label{eq:l07-dispatch}
    \end{equation}
    where $P(n)$ is the set of instructions that produce the operands of $n$
    (its closest preceding writers), and in which the execution unit is
    free:\tip{$P(n)$: for each register $n$ reads, its closest preceding
    writer (Lesson~6).}
    not held by another instruction $m$, which holds it in cycles $D_m$ to
    $W_m - 1$, and not taken by an older instruction that is also ready.
  \item[Write result] $W_n = D_n + k_n$, or later if the bus is taken in that
    cycle by a result with priority.
\end{description}
A unit or the bus may be taken by a \emph{later} instruction that became ready
earlier, so the rows cannot simply be completed in program order.\tip{Each
step waits for one resource: issue for a free station, dispatch for the
operands and a free unit, write result for the bus.}  Units and
the bus are assigned in cycle order across all instructions: in each cycle a
free unit takes the oldest of the instructions ready in that cycle, and the
bus the result with priority.  A unit never stays idle for an older
instruction that is not yet ready.

\begin{figure}[htbp]
  \centering
  % Timing of the running example: issue, waiting in the reservation station,
% execution, completed result waiting for the CDB, and write result.
\begin{tikzpicture}[x=0.52cm,y=0.5cm, font=\scriptsize\sffamily,
  c/.style={minimum width=0.52cm, minimum height=0.5cm, inner sep=0pt, draw=ln-rule},
  iss/.style={c, fill=white, draw=black},
  ex/.style={c, fill=ln-accent!25, draw=black},
  wt/.style={c, fill=ln-caution!20, draw=black},
  wr/.style={c, fill=ln-accent, draw=black, text=white, font=\scriptsize\bfseries\sffamily},
  q/.style={c, fill=ln-muted!30, draw=black}]
  \foreach \c in {1,...,19} { \node[text=ln-muted] at (\c,1) {\c}; }
  \node[anchor=east, text=ln-muted] at (0.4,1) {\iflangja{サイクル}{cycle}};
  \foreach \c in {1,...,19} { \draw[ln-rule, very thin] (\c-0.5,0.6) -- (\c-0.5,-6.5); }
  \draw[ln-rule, very thin] (19.5,0.6) -- (19.5,-6.5);
  % name, issue, first execution cycle, last execution cycle, write result
  \foreach \n/\op/\i/\ea/\eb/\w [count=\r from 0] in
    {1/MUL/1/2/5/6, 2/ADD/2/7/8/9, 3/SUB/3/4/5/7, 4/ADD/4/9/10/11,
     5/DIV/5/8/11/12, 6/MUL/7/12/15/16, 7/SUB/8/17/18/19} {
    \node[anchor=east] at (0.4,-\r) {I\n\enskip\texttt{\op}};
    \node[iss] at (\i,-\r) {I};
    \pgfmathtruncatemacro{\wa}{\i+1}
    \pgfmathtruncatemacro{\wb}{\ea-1}
    \ifnum\wa<\ea
      \draw[ln-muted, thick] (\wa-0.5,-\r) -- (\wb+0.5,-\r);
    \fi
    \foreach \e in {\ea,...,\eb} { \node[ex] at (\e,-\r) {}; }
    \node[wr] at (\w,-\r) {W};
  }
  % I3: result complete in cycle 5, bus taken by I1 in cycle 6
  \node[wt] at (6,-2) {};
  % I6: no free multiplier station in cycle 6
  \node[q] at (6,-5) {};
  % ---- legend
  \begin{scope}[shift={(0,-7.6)}]
    \node[iss] at (1,0) {I}; \node[anchor=west] at (1.5,0) {\iflangja{発行}{issue}};
    \draw[ln-muted, thick] (3.9,0) -- (4.9,0);
    \node[anchor=west] at (5,0) {\iflangja{RS で待機}{waiting in RS}};
    \node[ex] at (9.4,0) {}; \node[anchor=west] at (9.9,0) {\iflangja{実行}{execution}};
    \node[wr] at (13.3,0) {W}; \node[anchor=west] at (13.8,0) {\iflangja{結果書き込み}{write result}};
  \end{scope}
  \begin{scope}[shift={(0,-8.7)}]
    \node[wt] at (1,0) {};
    \node[anchor=west] at (1.5,0) {\iflangja{実行済み，CDB 待ち}{done, waiting for the CDB}};
    \node[q] at (9.4,0) {};
    \node[anchor=west] at (9.9,0) {\iflangja{空き RS なし（命令キュー）}{no free RS (in queue)}};
  \end{scope}
\end{tikzpicture}

  \caption{Timing of the running example.  Out-of-order execution lets I3
    run while I2 waits and I5 start while I4 waits; the delays of I3 (bus),
    I4 (adder) and I6 (reservation station) are those derived in
    \cref{ex:l07-timing}.}
  \label{fig:l07-timing}
\end{figure}

\begin{example}
  \label{ex:l07-timing}
  \Cref{tab:l07-timing} applies these rules to the running example, and
  \cref{fig:l07-timing} draws the result.  I3, I4 and I6 show every rule.
  I3 is dispatched in cycle~4, one cycle after its issue, and finishes
  executing in cycle~5; its write result would be cycle~6, but the bus then
  carries the result of I1, so $W_3 = 7$.  Whether the adder is free for I2 in
  cycle~7 thus depends on the row of the later I3.  I4 is ready from cycle~5,
  but the adder is occupied by I3, whose result waits for the bus, until its
  write result in cycle~7, and then by the older I2 until cycle~9, so
  $D_4 = 9$ and $W_4 = 11$.  I6 finds both multiplier
  stations busy until ML1 is freed by $W_1 = 6$, so $I_6 = 7$; its operands
  come from I2 and I4, and \cref{eq:l07-dispatch} gives
  $D_6 \ge 1 + \max(7, 9, 11) = 12$, which is also the cycle in which the
  multiplier is released by I5, so $D_6 = 12$ and $W_6 = 16$.  The table
  agrees with the full trace of \cref{tab:l07-trace}.\tip{With a single bus no
  two instructions broadcast in the same cycle: two equal entries in the $W_n$
  column mean that a bus conflict was missed.}
\end{example}

\begin{table}[htbp]
  \centering\small
  \caption{Issue, dispatch and write-result cycles of the running example.}
  \label{tab:l07-timing}
  \begin{tabular}{@{}llccc>{\raggedright\arraybackslash}p{0.44\linewidth}@{}}
    \toprule
    & Instruction & $I_n$ & $D_n$ & $W_n$ & What determines the timing \\
    \midrule
    I1 & \texttt{MUL R1, R2, R3} & 1 & 2  & 6  & ready at issue \\
    I2 & \texttt{ADD R4, R1, R5} & 2 & 7  & 9  & $W_1 = 6$; adder busy with I3 until $W_3 = 7$ \\
    I3 & \texttt{SUB R6, R2, R5} & 3 & 4  & 7  & ready at issue; bus taken in 6 \\
    I4 & \texttt{ADD R1, R5, R3} & 4 & 9  & 11 & adder: I3, then the older I2 \\
    I5 & \texttt{DIV R2, R6, R3} & 5 & 8  & 12 & $W_3 = 7$ \\
    I6 & \texttt{MUL R3, R4, R1} & 7 & 12 & 16 & no station until $W_1 = 6$;
                                               $W_4 = 11$; multiplier until $W_5 = 12$ \\
    I7 & \texttt{SUB R5, R3, R2} & 8 & 17 & 19 & $W_6 = 16$ \\
    \bottomrule
  \end{tabular}
\end{table}

\begin{keypoint}[Lesson summary]
  Tomasulo's algorithm combines register renaming with out-of-order
  execution, in which each instruction executes as soon as its operands are
  available.  Issue moves instructions in program order from the
  instruction queue into reservation stations, taking each operand either as
  a value from the register file or as a tag from the RAT, and points the RAT
  entry of the destination to the station.  A station captures values from
  the common data bus until its instruction is ready; dispatch then sends it
  to a free execution unit, choosing the oldest when several are ready.
  Write result broadcasts the value and tag, frees the station, and updates
  the register file and the RAT only if the result is not stale.  One bus
  serves one result per cycle, the slower unit having priority.  Loads and
  stores are performed in program order.  With one step per cycle, the issue,
  dispatch and write-result cycles of every instruction follow from its
  operands, the available stations and units, and the bus.  Tomasulo's
  algorithm does not support exceptions; the reorder buffer of the next lesson
  adds this.\margin{Issue, dispatch and write result stay exactly as
  described; Lesson~8 adds a fourth step, commit, which writes the register
  file in program order.}
\end{keypoint}

\end{document}
